% ===== PRÉAMBULE CANONIQUE — à coller VERBATIM en tête de CHAQUE .tex =====
\documentclass[11pt,a4paper]{article}
\usepackage[utf8]{inputenc}\usepackage[T1]{fontenc}\usepackage{lmodern}
\usepackage[french]{babel}
\usepackage{amsmath,amssymb,mathtools}\usepackage{icomma}
\usepackage{siunitx}\sisetup{locale=FR}  % \qty{}{}, \si{}, \ang{}
\usepackage{xcolor}\usepackage{colortbl}
\usepackage{tikz}\usepackage{pgfplots}\pgfplotsset{compat=1.18}
\usepackage[siunitx,europeanresistors]{circuitikz} % circuits (élec.)
\usepackage[most]{tcolorbox}\usepackage{enumitem}\usepackage{array,booktabs}
\usepackage[a4paper,margin=2.2cm]{geometry}
\usetikzlibrary{arrows.meta,calc,angles,quotes,positioning,patterns,%
  backgrounds,shapes.geometric,decorations.pathmorphing,decorations.markings,intersections}
\definecolor{primary}{RGB}{222,49,99}\definecolor{primarydark}{RGB}{150,22,55}
\definecolor{defcol}{RGB}{0,114,178}\definecolor{methcol}{RGB}{52,92,148}
\definecolor{excol}{RGB}{0,135,90}\definecolor{attncol}{RGB}{200,40,55}
\tcbset{boxbase/.style={enhanced,breakable,boxrule=0.9pt,arc=2.5pt,
  left=3mm,right=3mm,top=2.3mm,bottom=2.3mm,fonttitle=\bfseries,coltitle=white}}
% --- boîtes TUTORIEL (noms EXACTS, ne pas renommer) ---
\newtcolorbox{cle}[1][]{boxbase,colback=defcol!6,colframe=defcol,
  title={\textsc{À retenir}\ifx\relax#1\relax\relax\else\textmd{ : \itshape #1}\fi}}
\newtcolorbox{methode}[1][]{boxbase,colback=methcol!7,colframe=methcol,sharp corners,
  title={\textsc{Méthode}\ifx\relax#1\relax\relax\else\textmd{ --- \itshape #1}\fi}}
\newtcolorbox{exemplebox}[1][]{boxbase,colback=excol!5,colframe=excol,
  title={\textsc{Exemple}\ifx\relax#1\relax\relax\else\textmd{ : \itshape #1}\fi}}
\newtcolorbox{piege}[1][]{boxbase,colback=attncol!6,colframe=attncol,
  title={\textsc{Piège}\ifx\relax#1\relax\relax\else\textmd{ : \itshape #1}\fi}}
\newtcolorbox{remarque}[1][]{boxbase,colback=black!4,colframe=black!45,
  title={\textsc{Remarque}\ifx\relax#1\relax\relax\else\textmd{ : \itshape #1}\fi}}
% --- boîtes EXERCICE / CORRIGÉ (noms EXACTS, auto-comptées) ---
\newtcolorbox[auto counter]{exo}[1][]{boxbase,colback=primary!4,colframe=primary,
  title={\textsc{Exercice}~\thetcbcounter\ifx\relax#1\relax\relax\else\textmd{ --- \itshape #1}\fi}}
\newtcolorbox[auto counter]{corrige}[1][]{boxbase,colback=excol!4,colframe=excol,
  title={\textsc{Corrigé}~\thetcbcounter\ifx\relax#1\relax\relax\else\textmd{ --- \itshape #1}\fi}}
\newcommand{\vv}[1]{\vec{#1}}\newcommand{\ds}{\displaystyle}
\newcommand{\R}{\mathbb{R}}\newcommand{\N}{\mathbb{N}}\newcommand{\Z}{\mathbb{Z}}
\begin{document}
\sisetup{per-mode=symbol}

\begin{corrige}[vrai ou faux sur la nature de l'image]
\begin{enumerate}[label=\alph*)]
  \item \textbf{Vrai} : c'est le régime loupe ($0<u<f$) --- image virtuelle, droite, agrandie.
  \item \textbf{Faux} : pour $f<u<2f$, l'image réelle est \emph{agrandie} (principe du projecteur).
  \item \textbf{Faux} : une image virtuelle est \emph{toujours droite} ; « virtuelle et renversée » est
        impossible.
  \item \textbf{Vrai} : une divergente donne toujours, pour un objet réel, une image virtuelle,
        droite, réduite.
  \item \textbf{Faux} : une image réelle est \emph{toujours renversée}.
\end{enumerate}
Affirmations correctes : \textbf{a)} et \textbf{d)}.
\end{corrige}

\begin{corrige}[le rayon issu de la base de l'objet]
Tous les rayons issus d'un même point objet \textbf{convergent vers son image}. La base $A$ étant sur
l'axe, son image $A'$ est aussi sur l'axe : le rayon émergent doit donc passer par $A'$.
Le seul tracé qui passe par $A'$ est le \textbf{rayon 2}.\\
(1 = non dévié : faux ; 3 s'éloigne de l'axe ; 4 recroise l'axe trop tôt, avant $A'$.)
\end{corrige}

\begin{corrige}[l'objet glisse de $2F$ vers $F$]
En quittant $2F$ vers $F$, l'objet reste au-delà du foyer : l'image est réelle et grandit.
\begin{itemize}[label=--,leftmargin=1.4em,topsep=2pt,itemsep=1pt]
  \item a) \textbf{correcte} : l'image part de $2F'$ et file vers l'infini.
  \item b) \textbf{correcte} : elle reste réelle et renversée.
  \item c) fausse : l'image \emph{grandit} (de la même taille en $2F'$ jusqu'à très agrandie).
  \item d) fausse : l'image réelle se forme de l'\emph{autre} côté de la lentille.
\end{itemize}
Affirmations correctes : \textbf{a)} et \textbf{b)}.
\end{corrige}

\begin{corrige}[trois situations, trois instruments]
\begin{enumerate}[label=\alph*)]
  \item image \textbf{réelle, renversée, réduite} $\Rightarrow$ \textbf{appareil photo} (capteur).
  \item image \textbf{réelle, renversée, agrandie} $\Rightarrow$ \textbf{projecteur} (écran).
  \item image \textbf{virtuelle, droite, agrandie} $\Rightarrow$ \textbf{loupe} (\oe il).
\end{enumerate}
\end{corrige}

\end{document}
